\documentclass[10pt]{amsart}
\usepackage[a4paper,margin=23mm]{geometry}
\usepackage{amsmath,amssymb,amsthm,mathtools,hyperref,xfrac,enumitem}
\swapnumbers
\newtheorem{theo}[subsection]{Theorem}
\newtheorem{prop}[subsection]{Proposition}
\newtheorem{lemm}[subsection]{Lemma}
\newtheorem{coro}[subsection]{Corollary}
\theoremstyle{definition}
\newtheorem{defi}[subsection]{Definition}
\newtheorem{exem}[subsection]{Example}
\theoremstyle{remark}
\newtheorem{rema}[subsection]{Remark}
\emergencystretch=3em
\newcommand{\eg}{e.g.\ }\newcommand{\ie}{i.e.\ }\newcommand{\cf}{cf.\ }
\usepackage[scr=rsfs,cal=euler]{mathalfa}
\newenvironment{enumeratei}
{\bgroup\def\theenumi{\roman{enumi}}\def\theenumii{\arabic{enumii}}\begin{enumerate}}
{\end{enumerate}\egroup}
\newenvironment{enumeratea}
{\bgroup\def\theenumi{\alph{enumi}}\begin{enumerate}}
{\end{enumerate}\egroup}
\newcommand{\bbullet}{{\scriptscriptstyle\bullet}}
\newcommand{\cbbullet}{{\raisebox{1pt}{$\bbullet$}}}
\newcommand{\psfrac}[2]{\sfrac{(#1)}{#2}}
\newcommand{\Psfrac}[2]{(\sfrac{#1}{#2})}
\let\oldimplies\implies
\renewcommand\implies{\mathchoice{\oldimplies}{\Rightarrow}{\Rightarrow}{\Rightarrow}}

\makeatletter
\newcounter{paragraphe}
\renewcommand\theparagraphe{\thesubsubsection.\arabic{paragraphe}}
\makeatother

\hyphenation{retrac-tion defi-na-ble defect-less Berko-vich}

\usepackage{scalerel}
\usepackage{tikz-cd}


\newcommand{\opcit}{op.\@~cit.\@}

\newcommand{\abs}[1]{\mathopen|#1\mathclose|}

\newcommand{\an}{^{\mathrm{an}}}
\newcommand{\cl}{\mathrm{cl}}
\renewcommand{\d}{\mathrm d}
\newcommand{\val}{\mathrm{val}}
\renewcommand{\H}{\mathrm H}

\newcommand{\gm}{\bG_{\mathrm m}}
\newcommand{\tp}{\mathrm{tp\;}}
\newcommand{\acvsf}{\mathrm{ACV}^2\mathrm{F}}
\newcommand{\vsf}{\mathrm{V}^2\mathrm{F}}
\newcommand{\acvf}{\mathrm{ACVF}}
\newcommand{\DOAG}{\mathrm{DOAG}}
\newcommand{\RES}{\mathrm{RES}}
\newcommand{\res}{\mathrm{res}}
\newcommand{\VF}{\mathrm{VF}}

\newcommand{\bA}{\mathbb A}
\newcommand{\bG}{\mathbb G}
\renewcommand{\P}{\mathbb P}

\newcommand{\F}{\mathbb F}
\newcommand{\N}{\mathbb N}
\newcommand{\Q}{\mathbb Q}
\newcommand{\R}{\mathbb R}
\newcommand{\Z}{\mathbb Z}
\renewcommand{\SS}{\mathbb S}

\let\phi\varphi
\let\epsilon\varepsilon
\let\subset\subseteq

\renewcommand{\labelitemi}{$\bbullet$}
\renewcommand{\labelitemii}{$\diamond$}

\theoremstyle{plain}
\newtheorem{claim}[subsection]{Claim}
\newtheorem*{mtheo}{Main Theorem}

\theoremstyle{definition}
\newtheorem*{Examples*}{Examples}

\newcommand\stda[1]{{#1}^{\#}}

\title{Lattice generation of Lipschitz functions over divisible ordered abelian groups}
\author{Antoine Ducros; Ehud Hrushovski; Fran\c{c}ois Loeser; Jinhe Ye}
\date{}
\begin{document}
\maketitle
\noindent\textbf{Source.} Original Theorem 3.13 in Tropical functions on a skeleton, Journal de l'\'Ecole polytechnique -- Math\'ematiques 11 (2024), 613--654, DOI 10.5802/jep.261. The complete DOAG definitions, weak/lattice generation criteria, parameter descent, linear-cell separation and convex-subgroup induction are retained below. Ordinary cited special cell decomposition/model completeness facts remain precisely cited. The wider valued-field and Berkovich-skeleton results are outside this selected independent unit.
\setcounter{section}{2}
\section{Finite generation and Lipschitz functions in \texorpdfstring{$\DOAG$}{DOAG}}\label{section4}
In this section, we work in the theory of divisible ordered abelian groups which is denoted by $\DOAG$, and by definable we mean definable with parameters. We shall usually denote by
$\Gamma$ a model of $\DOAG$. We start with the definition of $w$-combination and $w$-generation.
\begin{defi}\label{def:t-gen}Let $X$ and $Y$ be definable topological spaces and $g$, $f_1,\ldots,f_n$ be definable continuous functions from $X$ to $Y$. We say $g$ is a \emph{$w$-combination} of $f_1,\ldots,f_n$ if for every $x\in X$, there is some $i\in \{1, \ldots, n\}$ such that $f_i(x)=g(x)$. Notationally, we use $[g=f_i]$ to denote the set $\{x\in X:g(x)=f_i(x)\}$. Hence, $g$ is a $w$-combination by $f_1,\ldots,f_n$ iff $X=\bigcup^n_{i=1}[g=f_i]$.
\end{defi}

In contrast, there is a stronger notion of combination that is very specific to $\DOAG$.
\begin{defi}\label{def:l-gen}Let $X$ be~a~definable topological space and let $g$ and $f_i$, $i\in I$, be~definable continuous functions $X\to \Gamma$. We say that $g$ is an \emph{$\ell$-combination} of the~$f_i$ if~$g$ lies in the $(\min,\max)$-lattice generated by $(f_i)_{i\in I}$. More explicitly, there are $f_1,\ldots,f_n$ in $(f_i)_{i\in I}$ such that $g$ is a function obtained by $f_1,\ldots,f_n$ and finitely many operations of $\min,\max$.
\end{defi}

We shall also use the following variants of $w$ and $\ell$-combination.
\begin{defi}\label{def:l-genvar}Let $X$ be a definable topological space and let $g$ and $f_i$
be definable continuous functions $X\to \Gamma$ for $i\in I$.
We say that $g$ is a \emph{$(w, +)$-combination} of the $f_i$ if there exist $h_1, \ldots, h_n$ in the abelian group generated by the functions $f_i$, $i \in I$ such that $g$ is a $w$-combination of the $h_i$.
We say that $g$ is an \emph{$(\ell, +)$-combination} of the $f_i$ if $g$ can be described by a formula involving only $+,-,\min$ and $\max$ and finitely many $f_i$.
\end{defi}

We say that a given set of functions containing the $f_i$ and stable under $w$\nobreakdash-combi\-nation is $w$-generated by the $f_i$ if it consists precisely of the set of all $w$-combinations of the $f_i$. We define $(w,+)$, $\ell$ and $(\ell,+)$-generation in an analogous way.

\begin{exem}\label{eg:min-k}
Let $X=\Gamma^n$ and $m_k:X\to \Gamma$ be the definable function which to $(x_1,\ldots,x_n)$ assigns the $k$-th smallest $x_i$. Clearly, $m_k$ is a $w$-combination of the coordinate functions $x_1,\ldots,x_n$. On the other hand, it is not hard to see that
\[
m_k(x)=\min_{\substack{U\subseteq \{1,\ldots,n\}\\|U|=k}} \max_{i\in U} x_i
\]
Hence the $m_k(x)$ are even $\ell$-combinations of $x_1,\ldots,x_n$.
\end{exem}
However, the two notions of combinations do not agree in general.
\begin{exem}\label{eg:dl}
Let $I$ be the interval $[0,\infty)\subseteq \Q$.
Let $D= I\times \{1,2\}\subseteq\Q^2$ and $f_1=0$,
$f_2=x_1$. Consider $g$ that is equal to
$f_i$ on $I\times \{i\}$ for $i=1,2$. Clearly
$g$ is a $w$-combination of $f_1$ and $f_2$.
However, we claim that $g$ is not an $(\ell, +)$-combination
of coordinate functions.
Indeed, if it were, then it would extend to a continuous
$\Q$\nobreakdash-definable function $g'$ on $\Q^2$. Let $\Gamma$ be a model of $\DOAG$
containing $\Q$ and in which there is some $c>n$ for all $n\in\N$. Since $\tp(1,c)=\tp(\alpha,c)$ for any $1>\alpha>0$ and $g(1,c)=c$, so $g'(\alpha,c)=c$. However $g'(0,c)=g(0,c)=0$, in contradiction with the continuity of $g'$. For a connected version of this example, replace $D$ by $D'=D\cup (\{0\}\times [1,2])$ and set $g=0$ on $\{0\}\times [1,2]$.
\end{exem}
This example suggests that interaction of the ambient space and the topology of~$D$ plays a role in distinguishing the two notions of combinations. To proceed towards a topological characterization for such properties, we need the following.
\begin{defi}\label{def:convex} Let $T$ be an o-minimal expansion of $\DOAG$ and $\Gamma\models T$ with $D\subseteq \Gamma^n$ definable. We say that $D$ is convex if for any $u$
and $v$ in $D$, $\psfrac{u+v}{2}\in D$.
\end{defi}
\begin{rema}\label{rmk:rcf}
When $T$ is an o-minimal expansion of the theory of real closed fields $\mathrm{RCF}$, this is equivalent to the usual definition of convexity for definable sets. For $u,v\in D$, let $L\subseteq [0,1]$ be $\{\alpha: \alpha u+(1-\alpha)v \in D\}$. By our notion of convexity, $L$~contains $\Z[1/2]\cap [0,1]$. By o-minimality, $L$ must be $[0,1]$ with at most finitely many points in $(0,1)$ removed. But removing any point from $(0,1)$ would lead to a violation of convexity.

Note further that for $D$ convex, working inside the smallest affine subspace containing $D$, we may assume that $\cl(\mathrm{int}(D))=\cl(D)$.
\end{rema}

Lastly, recall that for any definable subset $D$ of some $\Gamma^n$, a function $f:D\to \Gamma$ is called \emph{$\Q$-affine} if $f=\sum^n_{i=1} m_ix_i+c$ where $m_i\in \Q$ and $c\in \Gamma$. Such functions are the most basic definable continuous functions on $D$. We say $f$ is \emph{$\Z$-affine} if the $m_i$ are all in $\Z$.
\begin{prop}\label{prop:t-l-equiv}
Let $\Gamma$ be a divisible ordered abelian group and let $f_1,\ldots,f_m$ be $\Q$-affine functions on $\Gamma^n$. Let $D\subseteq \Gamma^n$ be definable and $g:D\to \Gamma$ be a continuous definable function. Assume that $g$ is a $w$-combination of $f_1,\ldots,f_m$. Then the following are equivalent:
\begin{enumerate}
\item\label{prop:t-l-equiv1} $g$ is an $\ell$-combination of $f_1,\ldots,f_m$.
\item\label{prop:t-l-equiv2} $g$ extends to a continuous definable function $g':\Gamma^n\to \Gamma$ that is a $w$\nobreakdash-combi\-na\-tion of $f_1,\ldots,f_m$.
\item\label{prop:t-l-equiv3} $g$ extends to a continuous definable function $g':D'\to \Gamma$ on some convex definable set $D'$ containing $D$ that is a $w$-combination by $f_1,\ldots,f_m$.
\item\label{prop:t-l-equiv4} For any $x,y\!\in\! D$, there is $i \!\in\! \{1,\ldots,m\}$ such that $f_i(x)\!\leq\! g(x)$ and $g(y)\!\leq\! f_i(y)$.
\item\label{prop:t-l-equiv5} For some collection $S$ of subsets of $\{1,\ldots,m\}$, $g=\min_{X\in S} \max_{i \in X} f_i$.
\end{enumerate}
\end{prop}

\begin{proof}
The implications $\eqref{prop:t-l-equiv5}\implies \eqref{prop:t-l-equiv1}\implies \eqref{prop:t-l-equiv2}\implies \eqref{prop:t-l-equiv3}$ are clear.

For $\eqref{prop:t-l-equiv3}\implies \eqref{prop:t-l-equiv4}$, by working in an elementary extension, we may assume that $\Gamma$ is a model of the theory of real closed fields $\mathrm{RCF}$. By Remark~\ref{rmk:rcf} and after replacing~$D$ by the convex set $D'$ in \eqref{prop:t-l-equiv3}, we may assume the line segment $[x,y]$ connecting $x,y$ is in $D$. Replace $g$ by $g'$ given by \eqref{prop:t-l-equiv3} as well. Let $I_j\subseteq [x,y]$ be $\{z: g(z)=f_j(z)\}$. By~continuity of $g$ and o-minimality, we know that the sets $I_j$ are finite unions of closed intervals and $\bigcup^m_{j=1} I_j=[x,y]$. Consider the canonical parameterization $h:[0,1]\to [x,y]$, $\alpha\mapsto \alpha y+(1-\alpha)x$, and let $f_i'=f_i\circ h$, $g'=g\circ h$ and $I'_j=h^{-1}(I_j)$. Since the functions $f_i$ are $\Q$-affine, the functions $f'_i$ are of the form $a_ix+b_i$ for some $a_i,b_i\in \Gamma$. Let $k$ be the $j$ such that $a_j$ is the greatest amongst all the $j$ such that $I'_j\neq \emptyset$. If~there are multiple such $j$, pick any. By induction, for $a$ to the right of $I'_k$, we have $g'(a)\leq f'_k(a)$. Similarly, for $a$ to the left of $I'_k$, we have $f'_k(a)\leq g'(a)$. In particular we have $f'_k(0)=f_k(x)\leq g'(0)=g(x)$ and $g'(1)=g(y)\leq f_k(y)=f'_k(1)$.

For $\eqref{prop:t-l-equiv4}\implies \eqref{prop:t-l-equiv5}$, consider $S$ to be the collection of subsets $X\subseteq \{1,\ldots,m\}$ such that $g\leq \max_{i\in X} f_i$ on the entire $D$.
Set $f:=\min_{X\in S}\max_{i\in X} f_i$.
We claim that $g= f$. Clearly $g\leq f$, so it suffices to show that $g\geq f$. For each $W\notin S$, there is some $y_W$ such that $g(y_W)>f_i(y_W)$ for
every $i \in W$. By \eqref{prop:t-l-equiv4}, for each $x\in D$, there is $i^x_{W}$ such that $f_{i^x_{W}}(x)\leq g(x)$ and $f_{i^x_W}(y_W)\geq g(y_W)$. Note that $i^x_W\notin W$. Let $X=\{i^x_W:W\notin S\}$. We have that $X\in S$ because otherwise, $i^x_X\in X$. For this $x$, we have that $\max_{i\in X} f_i(x)\geq g(x)$ and $f_i(x)\leq g(x)$ for any $i\in X$, hence $f(x)\leq \max_{i\in X} f_i(x)=g(x)$.
\end{proof}

\begin{coro}\label{cor:cordinate}
Let $D\subseteq \Gamma^n$ be a definable convex set. The set of definable continuous functions from $D$
to $\Gamma$ is $(\ell,+)$-generated by the constants and all rational multiples of coordinate functions.
\end{coro}
\begin{proof}
By quantifier elimination, we can find $\Q$-affine functions $f_1,\ldots,f_n$ such that $g$ is a $w$-combination of
$f_1,\ldots,f_n$. By Proposition~\ref{prop:t-l-equiv}, we have that $g$ is in fact an $\ell$-combination of $f_1,\ldots,f_n$.
\end{proof}
Proposition~\ref{prop:t-l-equiv} suggests that the agreement of $w$-combination and $\ell$-combination is related to the existence of continuous extensions to an ambient convex space. This motivates the following definition.
\begin{defi}\label{def:Lipschitz}
For a tuple $x\in \Gamma^n$, define $|x|=\max^n_{i=1}|x_i|$. Let $D\subseteq\Gamma^n$ and $f:D\to \Gamma$ a definable function. We say $f$ is \emph{Lipschitz} if there is some $M\in \N$ such that $|f(x)-f(y)|\leq M|x-y|$. \end{defi}

Note that Lipschitz functions are automatically continuous and clearly the class of Lipschitz functions depends on the embedding of $D$ in $\Gamma^n$.
Our purpose is now to investigate
Lipschitz definable functions on closed definable sets; a first step
will consist in reducing to the definably compact case,
by using the two following lemmas.

\begin{lemm}\label{lem:lipschitz-closure}
Let $\Gamma$ be a model of $\DOAG$, let $D$
be a subset of $\Gamma^n$ definable over some set
$A$ of parameters, and let
$f\colon D\to \Gamma$ be a Lipschitz
$A$-definable map. Let $(f_i)$ be a finite family
of $\Q$-affine $A$-definable functions
such that $f$ is a $w$-combination
of the $f_i|_D$. Then $f$ admits a unique
continuous extension $\overline f$ to $\cl (D)$,
the set $\cl (D)$ and the function $\overline f$
are $A$-definable, and $\overline f$ is Lipschitz
and if a $w$-combination of
the $f_i|_{\cl (D)}$.
\end{lemm}

\begin{proof}
The uniqueness of $\overline f$ is clear, as well
as the $A$-definability of $\cl (D)$ and $\overline f$
if the latter exists, as one sees by
using the definition of the closure and of the limit
(with $\epsilon$ and $\delta$\ldots). The same reasoning
also shows that the set of points of $\cl (D) \setminus D$ at which $f$
admits a limit is $A$-definable.
Moreover if $\overline f$ exists it inherits obviously the Lipschitz
property of~$f$, and it is also $w$-generated by the (restrictions of)
the~$f_i$: indeed, the subset of $\cl (D)$ consisting of points $x$
such that there is some $i$ with $f(x)=f_i(x)$ is closed and contains $D$, thus is the whole of $\cl (D)$.

It thus remains to show the existence of $\overline f$,
and this can be done after enlarging the model $\Gamma$.
We can thus suppose that it is equal to the additive group of some
real closed field.
Let $x$ be a point of $\cl (D) \setminus D$.
There exists a half-line $L$ emanating from
$x$ such that $(x,y)\subset D$ for some $y$;
taking $y$ close enough to $x$ we can assume that $f=f_j$ on
$(x,y)$ for some $j$. Then the limit of~$f$
at $x$ along the direction of $L$ exists and is equal to $f_j(x)$.
The Lipschitz property then ensures
that this limit does not depend on $L$,
let us
denote it by $\overline f(x)$. Since $D$
is defined by affine inequalities, there is a positive
$\gamma\in \Gamma$ such that for every $y$
in $\Gamma^n$ with $\|x-y\|<\gamma$
(say for the Euclidean norm) then either $(x,y)\subset D$
or $(x,y)\cap D=\emptyset$. Thus if $y$ is a point of $D$
with $\|x-y\|<\gamma$ then
$\abs{f(y)-\overline f(x)}\leq N\|x-y\|$
where $N$ is an upper bound for the slopes of the $f_i$.
So $f(y)$ tends to $\overline f(x)$ when the point $y$
of $D$ tends to $x$.
\end{proof}

\begin{lemm} \label{lem:comp-l-gen}
Let $M$ be either $\{0\}$ or a model of $\DOAG$, let
$\Gamma$ be a model of $\DOAG$
containing $M$, and let $\rho$
be an element of $\Gamma$
with $\rho >M$.
Let $Z\subseteq \Gamma^n$ be
an $M$-definable subset. Let $x_1,\ldots,x_n:Z \to \Gamma$ denote the coordinate functions of $Z$ and let $h:Z\to \Gamma$ be an
$M$-definable function.

Assume that there
exists a term $t$
in
$\{+,-,\max,\min\}$ and
$\gamma=(\gamma_1,\ldots,\gamma_l)$
in $\Gamma^\ell$
such that $h|_{Z_\rho}=t(x_1,\ldots,x_n,\gamma)|_{Z_\rho}$,
where $Z_\rho=Z\cap [-\rho,\rho]^n$.

Then there is a term $t'$ in $\{+,-,\max,\min\}$ and
a finite tuple $\beta$ of elements of~$M$ such that $h=t'(x_1,\ldots,x_n,\beta)$.
\end{lemm}

\begin{proof}
Assume first that $M$ is
a model of $\DOAG$. By our assumption,
there exists a term $t$ in $\{+,-,\max,\min\}$ and
a tuple
$\gamma=(\gamma_1,\ldots,\gamma_l)\in \Gamma^\ell$
such that $h|_{Z_\rho}=t(x_1,\ldots,x_n,\gamma)|_{Z_\rho}$.
By model-completeness of $\DOAG$, the $\gamma_i$ can be chosen in
$M\oplus \Q\cdot \rho$.
Thus there is $m>0$ such that for each $i$, there exist integers $k_i$ and $\beta_i\in M$
with $\gamma_i=\Psfrac{k_i}{m}\rho+\beta_i$. Let $\nu$ denote $\rho/m$.
We have
\[
h|_{Z_{m\nu}}=t(x_1,\ldots,x_n,(k_i\nu+\beta_i))|_{Z_{m\nu}}.
\]
Viewing the above as a first-order formula with constants in
the model $M$ and a variable for $\nu$, using o-minimality
and model-completeness of~$M$, we have some $\nu_0\in M_{>0}$ such that for any $\nu'>\nu_0$, the following holds in $M$:
\[
h|_{Z_{m\nu'}}=t(x_1,\ldots,x_n,(k_i\nu'+\beta_i))|_{Z_{m\nu'}}.
\]
Take $\nu(x)=\max \{|x_1|,\ldots,|x_n|,2\nu_0\}$ and
\[
t'(x_1,\ldots,x_n,\beta)=t(x_1,\ldots,x_n,(k_i\nu(x)+\beta_i)).
\]
We then have
\[
h=t'(x_1,\ldots,x_n,\beta)
\]
by construction, which ends the proof when $M\neq \{0\}$.

If $M=\{0\}$,
set $\Gamma'=\Gamma\oplus \Q\cdot \delta$ where $\delta$ is positive
and infinitesimal with respect to~$\Gamma$, set $M'=\Q\cdot \delta$ and
let us denote
by $Z'$ and $h'$ the objects
deduced from $Z$ and $h$ by base-change to $\Gamma'$.
Applying the above yields
a
term
$\theta$ in $\{+,-,\max,\min\}$
and a tuple $\beta$
of elements of $\Q\cdot \delta$ such that $h'=\theta(x_1,\ldots,x_n,\beta)$.
By reducing modulo the convex subgroup $\Q\cdot \delta$ of $\Gamma'$
we see that $h=\theta(x_1,\ldots,x_n,0)$.
\end{proof}

We can now state the main result of this section.

\begin{theo}\label{thm:lip-t-gen}
Let $M\models \DOAG$ or $M=\{0\}$ and
let $\Gamma$ be a model of
$\DOAG$ containing $M$.
Let $D\subseteq \Gamma^m$ be an $M$-definable set. Let $g:D\to \Gamma$ be a Lipschitz definable function over $M$. Let $f_1,\ldots,f_n$ be $\Q$-affine functions over $M$ such that $g$ is a $w$-combination of $f_1,\ldots,f_n$. Then $g$ is an $(\ell, +)$-combination of the $f_i$,
the constant $M$-valued functions and the coordinate functions.
\end{theo}

Before proving this result we will need some preliminaries on cell decomposition in $\DOAG$.

\subsection{Cell decomposition}
Fix a model $\Gamma$ of
$\DOAG$.
We shall use the notion
of special linear decompositions from~\cite{elef}. In~\cite{elef}, Eleftheriou defines the notion of linear decomposition, which is a cell decomposition using only
graphs of $\Q$-affine functions instead of general piecewise $\Q$-affine functions. In fact we will need only to consider bounded linear cells in $\Gamma^n$. They are defined by induction on $n$.
In $\Gamma^0$ the origin is a bounded linear cell. If $C$ is a bounded linear cell in $\Gamma^{n-1}$, $f$ and $g$ are $\Q$\nobreakdash-affine functions on $\Gamma^{n-1}$, with $f < g$ on $C$, the relative interval
\[
(f < g)_C = \{(x', y) \in C \times \Gamma; f (x') < y < g(x')\}
\]
and the graph $\Gamma (f)_C = \{(x', y) \in C \times \Gamma; f (x') = y\}$ are bounded linear cells in $\Gamma^n$.
If $Y$ is a bounded definable set in $\Gamma^n$, a linear decomposition of $Y$ is a partition of $Y$ into (finitely many) bounded linear cells.

We denote by $\pi: \Gamma^n \to \Gamma^{n-1}$ the projection to the $n-1$ first coordinates.
A~special linear decomposition of a bounded definable set $Y \subseteq \Gamma^n$ is defined recursively in~\cite{elef} as follows. When $n = 1$ any cell decomposition of $Y$ is special.
If $n > 1$, a~linear decomposition $\mathcal{C}$ of $Y$ is special if the following conditions are satisfied:\enlargethispage{\baselineskip}%
\begin{enumerate}
\item $\pi (\mathcal{C})$ is a special linear decomposition of $\pi (Y)$.
\item For every pair of cells $\Gamma (f)_S$ and $\Gamma (g)_T$ in $C$ with $S$ in the closure of $T$,
$f|_S <g|_S$ or $f|_S >g|_S$ or $f|_S =g|_S$.
\item\label{item3} For every pair of cells $(f<g)_T$ and $X$ in $C$, where $X = \Gamma (h)_S$, $(h,k)_S$ or $(k,h)_S$,
there is no $c \in \cl(S) \cap \cl(T)$ such that $f(c) < h(c) < g(c)$.
\end{enumerate}

An important property of special linear decompositions is that if $D$ and $E$ are two cells in such a decomposition such that
$D \cap \cl (E)$ is non-empty then $D \subseteq \cl (E)$ \cite[Fact 2.3]{elef}.
By~\cite[Fact 2.2]{elef} special linear decompositions of $Y$ always exist.

Note that closures of cells have a simple description:
the closure of $(f < g)_C$ is equal to $(f \leq g)_{\cl (C)} = \{(x', y) \in \cl (C) \times \Gamma; f (x') \leq y \leq g(x')\}$ and the closure of
$\Gamma (f)_C$, is $\Gamma (f)_{\cl (C)}$.
In particular, if $C$ is a cell, $\pi (\cl (C)) = \cl (\pi (C))$.

\begin{lemm}\label{lemm:cell}Fix a special linear cell decomposition of a closed bounded definable subset of $\Gamma^n$ and let $C_1$ and $C_2$ be two cells. Set $D_1 = \cl (C_1)$ and
$D_2 = \cl (C_2)$. Assume that $D_1 \cap D_2$ is non-empty. Then there exists a cell $C$ such that
$D_1 \cap D_2 = \cl (C)$.
\end{lemm}

\begin{proof}We proceed by induction on $n$.
The case $n=0$ is clear. If $n >0$, we have
that $\pi (D_1) \cap \pi (D_2) = \cl (C')$ for some cell $C'$ of the projection of the decomposition.
Since for $i = 1, 2$, $D_i\cap \pi^{-1} (C')$ is either of the form
$(f_i \leq g_i)_{C'}$ or $\Gamma (f_i)_{C'}$, it follows from condition \eqref{item3} of being a special linear decomposition that either
$D_1 \cap D_2 \cap \pi^{-1} (C')
= (f_1\leq g_1)_{C'}
= (f_ 2\leq g_2)_{C'}$ or
$D_1 \cap D_2 \cap \pi^{-1} (C')
= \Gamma (f_1)_{C'} = \Gamma (f_2)_{C'}$, from which the statement follows.
\end{proof}

We shall also need the following statement.

\begin{lemm}\label{lemm:ineg}Let $D$ be a closed bounded definable subset of
$\Gamma^n$. Assume $D$ is convex. Let $h$ be a $\Q$-affine function on $\Gamma^n$ such $h \geq 0$ on $D$.
Let $D_0$ be the zero locus of $h$ in $D$. We assume that $D_0$ is non-empty and $D_0 \neq D$.
Let $f$ be a $\Q$-affine function on $\Gamma^n$ which vanishes on $D_0$. Then there exists a positive integer
$M$ such that, for every $x \in D$, $\vert f (x) \vert \leq M h (x)$.
\end{lemm}

\begin{proof}Let $\mathcal{D}$ be a linear decomposition of $D$.
We consider the set $\mathcal{F}$ of all sets $F$ of the form $F = \cl (C)$, with $C$ a 1-dimensional cell in $\mathcal{D}$,
that intersect the hyperplane $h = 0$ and are not contained in $h=0$.
For such an $F$ we denote by $p_F$ its intersection point with $h = 0$.
There exists a positive integer $M_F$ such that
$\vert f (x) \vert \leq M_F h (x)$ on~$F$. Indeed,
the restrictions of both $h$ and $g$ to the line segment $F$ are linear functions on~$F$ vanishing at the endpoint $p_F$ of $F$ and the restriction of $h$ is not identically zero, which yields the existence of
some $M_F$.
In fact the inequality $\vert f (x) \vert \leq M_F h (x)$ holds on the whole half-line $L_F$ containing $F$ with origin $p_F$.
Take $M = \max_{\mathcal{F}} (M_F)$. Now consider $R$ a $\mathrm{RCF}$-expansion of $\Gamma$. Let $Y$ be the convex hull of the half-lines $L_F$ in that expansion. We have $\vert f (x) \vert \leq M h (x)$ on $Y$.
But $Y$ contains $D (R)$, since if $P$ is a convex definably compact polyhedron of $R^n$, and if $F$ is a face of $P$ of any dimension, then the convex hull of all half-lines directed by $1$-faces intersecting $F$ contains $P$,
hence the result, taking $P = D$ and $F = D_0$.
\end{proof}

The following statement about separation by hyperplanes will play a key role in our proof of Theorem \ref{thm:lip-t-gen}.

\begin{prop}\label{prop:cell}Fix a special linear cell decomposition of a closed bounded definable subset of $\Gamma^n$ and let $C_1$ and $C_2$ be two cells.
Assume $C_1 \neq C_2$. Set $D_1 = \cl (C_1)$ and
$D_2 = \cl (C_2)$. Then there exists a $\Z$-affine function $h$ such that $h \geq 0$ on $D_1$,
$h \leq 0$ on $D_2$, and
the hyperplane $H =h^{-1}(0)$ satisfies $D_1\cap D_2 =D_1 \cap H=D_2\cap H$.
\end{prop}

\begin{proof}We shall proceed by induction on $n$, the case $n=1$ being clear.
If $\pi (C_1) = \pi (C_2)$, then the statement is clear. Indeed, for each $i = 1, 2$, we have
$C_i = (f_i < g_i)_S$ or $C_i = \Gamma(f_i)_S$. In the second case we set $g_i = f_i$.
We may assume that $C_1$ is above~$C_2$. The graph of the average of $f_1$ and $g_2$ provides the required hyperplane.

Thus we will assume from now on that $\pi (C_1) \neq \pi (C_2)$.
We set $C'_i = \pi (C_i)$ for $i = 1, 2$.
By Lemma \ref{lemm:cell}, if $D_1 \cap D_2$ is non-empty, there exists a cell $C$ such that
$D_1 \cap D_2 = \cl (C)$.

\subsubsection*{Case 1: $D_1 \cap D_2$ is non-empty and $C$ is of the form $(f < g)_S$}

In this case, for $i = 1, 2$, $C_i$ is necessarily of the form
$(f_i < g_i)_{C'_i}$ where $f_i$ and $g_i$ are $\Q$-affine functions coinciding with $f$ and $g$ on $S$, since we are working with a special linear cell decomposition.
Furthermore $D_i = (f_i \leq g_i)_{\cl (C'_i)}$ and we have $f_1 = f_2$ and $g_1 = g_2$ on $\cl (C'_1) \cap \cl (C'_2)$.
It follows that $D_1 \cap D_2 =
(f_i \leq g_i)_{\cl (C'_1) \cap \cl (C'_2)}$, for $i = 1,2$.
By the induction hypothesis, there exists an hyperplane $h'$ in $\Gamma^{n-1}$ given by a
$\Z$-affine equation satisfying
the conditions of Proposition \ref {prop:cell} relatively to
$\cl (C'_1)$ and $\cl (C'_2)$. Consider the vertical hyperplane $H$ above $h'$ (the hyperplane defined by the same equation in $\Gamma^n$).
It follows from our description of $D_1 \cap D_2$ that $H$ satisfies the required conditions.

\subsubsection*{Case 2: $D_1 \cap D_2$ is non-empty and $C$ is of the form $\Gamma (f)_S$}

By the induction hypothesis,
there exists an hyperplane $h'$ given by an equation $h' (x')=0$ in $\Gamma^{n-1}$, with $h'$ a $\Z$-affine function, $x'=(x_1,\ldots, x_{n-1})$
fulfilling
the conditions of Proposition \ref {prop:cell} relatively to
$\cl (C'_1)$ and $\cl (C'_2)$.
In particular $h' \geq 0 $ on
$\pi (D_1)$ and $h' \leq 0$ on $\pi (D_2)$.
We denote by $H$ the hyperplane with equation $h'(x')=0$ in $\Gamma^{n}$.\enlargethispage{\baselineskip}%

The set $C_1$ is of the form $(f_1 < g_1)_{C'_1}$ or $\Gamma (f_1)_{C'_1}$. In the second case we set $g_1 = f_1$.
Similarly $C_2$ is of the form $(f_2 <g_2)_{C'_2}$ or $\Gamma (f_2)_{C'_2}$ and in the second case we set $g_2 = f_2$.
Set $C' = \pi (C)$.
Since our linear decomposition is special,
we~have that $f|_{C'}$ is equal to $f_1|_{C'}$ or $g_1|_{C'}$. Without loss of generality we may assume that $f|_{C'} = g_1|_{C'}$. It follows that $f|_{C'}=f_2|_{C'}$ by the case assumption and the fact our decomposition is special.
Let $X$ be the graph of $g_1$ over $\cl (C'_1)$. The function $x_n -\nobreak f (x')$ is identically zero on
$H \cap X$, hence by Lemma \ref{lemm:ineg}, there exists a positive integer~$M$ such that
$\vert x_n - f (x') \vert \leq M h' (x')$ on $X$. After increasing $M$ we may assume the inequality is strict when $h' (x') \neq 0$.
It follows that the
the hyperplane $H_M$ with equation \hbox{$Mh'(x') - (x_n - f (x')) = 0$} lies above the set $D_1$
and strictly above $D_1 \setminus\nobreak H$. Using the same argument for $D_2$, we get that after possibly increasing $M$
the hyperplane~$H_M$ lies under the set~$D_2$ and strictly under $D_2 \setminus H$.
Let us check that~$H_M$ satisfies the required conditions.
Indeed, a point $x = (x', x_n)$ lies in $D_1 \cap H_M$ if and only if
$x' \in \pi (D_1)$, $x \in H_M$, and $f_1 (x') \leq x_n \leq g_1 (x')$.
But if $x \in D_1 \cap H_M$ we must have $h'(x') = 0$.
Thus
$x$ lies in $D_1 \cap H_N$ if and only if
$x' \in \pi (D_1)$, $x \in H_M$, $h'(x')= 0$ and $x_n = f(x')$, from which the equality
$D_1 \cap H_N = D_1 \cap D_2$ follows, and one gets similarly that
$D_2 \cap H_N = D_1 \cap D_2$.

\subsubsection*{Case 3: $D_1 \cap D_2$ is empty}

If $\pi (D_1) \cap \pi (D_2) = \emptyset$ then by the induction hypothesis there exists an hyperplane
$h'$ in $\Gamma^{n -1}$ satisfying the required conditions for $\pi (D_1)$ and $\pi (D_2)$ and $\pi^{-1} (h')$ will do the job.
Thus we may assume that $\pi (D_1) \neq \pi (D_2)$ and $\pi (D_1) \cap \pi (D_2) \neq \emptyset$.
We choose an
hyperplane
$h'$ in $\Gamma^{n -1}$ with equation $h'(x')=0$ satisfying the required conditions for $\pi (D_1)$ and $\pi (D_2)$.
We may assume $h' \geq 0$ on~$D_1$ and $h' \leq 0$ on $D_2$.
As in Case 2, $C_1$ is of the form $(f_1 <g_1)_{C'_1}$ or $\Gamma (f_1)_{C'_1}$. In the second case we set $g_1 = f_1$.
Similarly for $C_2$.

Set $D'_1= D_1\cap H$ and $D'_2 = D_2 \cap H$.
We have $D'_1 \cap D'_2 = \emptyset$.
We may assume that $f_2 > g_1$ over $\pi (D_1) \cap \pi (D_2)$.
Note that if we intersect the cells of a special linear cell decomposition of some bounded set $W$ with $H$ we get a special linear cell decomposition of $W \cap H$.
Thus we can apply the induction hypothesis to~$D'_1$ and~$D'_2$, and there exists a $\Z$-affine function $f$ on $\Gamma^n$
such that $f >0$ on $D'_1$ and $f <0$ on~$D'_2$.
We claim that for $M$ a large enough integer the hyperplane
$M h' +f = 0$ will separate~$D_1$ and~$D_2$.

To prove this we proceed similarly as in the proof of Lemma \ref{lemm:ineg}.
We consider the set $\mathcal{F}$ of all sets $F$ of the form $F = \cl (C)$, with $C$ a 1-dimensional cell contained in~$D_1$,
that intersect $H$ and are not contained in $H$.
For such an $F$ denote by
$p_F$ the intersection point of $F$ with $H$.
The restriction of~$f$ to
$F$ can be written as $f (p_F) + \ell_F$ with $\ell_F$ a $\Q$-linear function on $F$.
Since $h'$ is strictly positive on $F$ outside $p_F$,
there exists a positive integer $M_F$ such that $M_F h' + \ell_F \geq 0 $ on $F$.
This still holds on the whole half-line $L_F$ containing $F$ with origin $p_F$.
Since $f (p_F) > 0$ by assumption, we~get that $M_F h' + f > 0$ on $L_F$.
Take $M_1 = \max_{\mathcal{F}} (M_F)$.
Proceeding as in the proof of Lemma \ref{lemm:ineg} we deduce that
$M_1 h' + f >0$ on $D_1$.
One proves similarly the existence of $M_2$ such that
$M_2 h' + f <0$ on $D_2$. Thus we can take $M = \max (M_1, M_2)$.
\end{proof}

\begin{proof}[Proof of Theorem \ref{thm:lip-t-gen}]
By Lemma \ref{lem:lipschitz-closure}
we can assume that $D$ is closed.
We may then enlarge the model $\Gamma$
and assume that it contains some
$\rho$ with $\rho>M$. Let
$D_\rho$ denote the
intersection of $D$ with $[-\rho,\rho]^m$.
This is a definably compact subset of $\Gamma^m$
which is definable over $M_\rho:=M\oplus \Q\cdot \rho$.
If we prove that $g|_{D_\rho}$
is an $(\ell,+)$-combination of the $f_i$, the coordinate
functions and some constant functions with values in $M_\rho$,
Lemma \ref{lem:comp-l-gen}
above will allow us to conclude that $g$
is an $(\ell,+)$-combination of the $f_i$, the coordinate
functions and some constant functions with values in $M$.
We thus may and do assume that $D$ is definably compact.
By considering a
submodel of~$M$ over which everything is defined, we reduce to
the case
where
$M$ has exactly $r$ non-trivial convex subgroups,
and we proceed by induction
on $r$. The case $r=0$
is obvious since the definably compact set $D$ is then either empty
or equal to
$\{0\}$. Assume now
that $r>0$ and that the result holds true
for smaller values of $r$.\enlargethispage{.5\baselineskip}%

Let $M_0$ be the smallest non-trivial convex subgroup of~$M$, and $\overline{M}=M/M_0$ be the quotient.
We are first going to explain why we can assume that $g(D(M))
\subseteq M_0$; this is tautological if $M_0=M$, so we assume
(just for this reduction step) that $M_0\neq M$.
In this case $\overline M$
is a model of $\DOAG$ with $r-1$ non-trivial convex subgroups, and the natural map carrying $M$ to $\overline{M}$ induces a map that carries $D(M)$ to a definably compact
definable subset $\overline{D}$ of $\overline{M}^m$ (see~\cite[Th.\,4.1.1]{BP} for example). Furthermore, since $g$ is Lipschitz, it
descends to a definable function $\overline{g}:\overline{D}(\overline{M})\to \overline{M}$, which is Lipshitz as well
and is a $(w,+)$-combinations of the $\overline{f_i}$.
By the induction hypothesis, we then know that $\overline{g}$ is of the form $\tau(\overline{f_1},\ldots,\overline{f_n})$, where $\tau$ is a term involving constants, projections and $+,-,\min,\max$ only. Replacing $g$ by $g-\tau(f_1,\ldots,f_n)$, we may assume that $g(D)\subseteq M_0$, as announced.

By~\cite[Fact 2.2]{elef} there exists a special linear decomposition $\mathcal{D}$ of $D$ such that~$g$ is $\Q$-affine on each cell.
Clearly $D$ is covered by the closed sets $D_i = \cl (C_i)$, for~$C_i$ in~$\mathcal{D}$.
In fact if one considers the set
$\mathcal{D}'$ of all $C\in \mathcal{D}$ such that, for any $C'\neq C$, $C$~is not contained in the closure of $C'$, it follows from
\cite[Fact 2.3]{elef} that the closed sets $D_i = \cl (C_i)$ for
$C_i$ in $\mathcal{D}'$ already cover $D$, but we will not use this. Sets of the form $\cl (C)$ with $C \in \mathcal{D}$ will be referred to as closed cells.

We will now use the separating hyperplanes provided by Proposition \ref{prop:cell}
to build affine functions that will appear in the $(\ell,+)$-combination we are seeking for describing~$g$. For this purpose,
the inclusion~$g(D(M))\subseteq M_0$ will be crucial.

\begin{claim}\label{claim:separation}
Let $C'$ and $C''$ be any two distinct cells in $\mathcal{D}$. Set $D' = \cl (C')$ and $D'' = \cl (C'')$.
There exists a function $f_{D',D''}$ in the group generated by $f_1,\ldots,f_n$, the constant functions and the coordinate functions such that\vspace*{-3pt}
\begin{equation}\label{eqn:separation}\tag{$*$}
g|_{D'}\leq f_{D',D''}|_{D'} \text{ and } f_{D',D''}|_{D''}\leq g|_{D''}.
\end{equation}
\end{claim}

\begin{proof}[Proof of the claim]
By Proposition~\ref{prop:cell} there exists a $\Z$-affine function $h$ such that the hyperplane $H =h^{-1}(0)$ satisfies $D' \cap D'' =D' \cap H=D''\cap H$, $h \geq 0$ on $D'$ and
$h \leq 0$ on $D''$.

If $D'\cap D''=\emptyset$, using definable compactness of $D'$ and $D''$, we get that there exists $a \in M_0$ such that $h|_{D''}<-a<0<a<h|_{D'}$.
Moreover, by our assumption that $g(D(M)))\subseteq M_0$, there is $b\in M_0$ such that $g(D (M)) \subseteq (-b, b)$.
For any positive integer $m$ we have $mh - g > ma - b$ on $D'$ and $mh - g < -ma + b$ on $D''$.
Since $M_0$ is archimedean, for $m$ large enough, we have $ma > b$, hence
condition~(\ref{eqn:separation}) is satisfied for $f_{D',D''} = mh$.

If $D'\cap D''\neq \emptyset$, take $c\in D'\cap D''$ and let $G$ be the $\Q$-affine function such that $g=G$ on $D'$.
Replacing $g$ by $g-G$, we may assume that $g=0$ on $D'$. Translating our entire set by $c$, we may assume that $c$ is the origin. Thus $g(0)=0$
and $g$ is actually the restriction of a $\Q$-linear function on $D''$.
On $D'$, for any positive integer $m$, we have $mh\geq0=g$.
For any $b\in D''$, if $h(b)=0$, then $b\in D''\cap H=D'\cap H$, hence $g(b)=0$.
Thus, by Lemma \ref{lemm:ineg}, there exists a positive integer $m$ such that
$-g \leq -mh$ on $D''$. For such an integer $m$, we have $g \leq mh$ on $D'$ and $g \geq mh$ on $D''$.
\end{proof}

We can now conclude the proof of Theorem \ref{thm:lip-t-gen}. Note that $g$ is a $w$-combination of the functions $f_i$; it is thus a fortiori a $w$-combination of
the set of functions obtained by adding all the functions $f_{D',D''}$ from Claim \ref{claim:separation} to the functions $f_i$.
Take~$x$ and~ $y$ in $D$. If they belong to the same closed cell $D' = \cl(C')$, then $g (x) = f_i (x)$
and $g (y) = f_i (y)$ for some $i$ and condition \eqref{prop:t-l-equiv4} in Proposition~\ref{prop:t-l-equiv} is satisfied.
If they belong to two distinct closed cells $D'$ and $D''$,
then $g (x) \leq f_{D',D''} (x)$ and $ f_{D',D''} (y) \leq g (y)$ by Claim \ref{claim:separation}.
Thus, by the implication $\eqref{prop:t-l-equiv4}\implies\eqref{prop:t-l-equiv1}$ in
Proposition~\ref{prop:t-l-equiv}, we obtain that~$g$ is an $\ell$-combination of
the functions $f_i$ and $f_{D',D''}$, which concludes the proof.
\end{proof}
The proof of Claim~\ref{claim:separation} actually yields the following convenient way to check if a given function is Lipschitz on a definably compact set.

\begin{coro}
Let $D'$ and $D''$ be two definably compact convex sets such that $D'\cap D''=D'\cap H_a=D''\cap H_a\neq \emptyset$ with $H_a$ an hyperplane defined by a $\Z$-affine function. Assume further that $g$ a continuous function that is affine on $D'$ and $D''$ respectively, then $g$ is Lipschitz on $D'\cup D''$.
\end{coro}

\begin{rema}
Note that one can have definably compact versions of Example~\ref{eg:dl} by replacing $[0,\infty)$ with $[0,c]$ for some $c>n\cdot 1$ for all $n\in \N$. However, the function~$g$ there is not Lipschitz because $|(0,c)-(1,c)|=1$ and $|g(0,c)-g(1,c)|=c$.
\end{rema}

\begin{rema}In the case of homogeneous linear equations, with no parameters, equivalence of $\ell$-combination and $w$-combination goes back to work of Beynon \cite{beynon}, see also \cite[\S5.2]{glass} and \cite{Ovchinnikov} for related results. In 2011, as a student, Daniel Lowengrub rediscovered and partially generalized Beynon's results. He also gave Exam\-ple \ref{eg:dl} showing that they do not hold over non-archimedean parameters. Here we fully generalized them, after replacing continuity by a Lipschitz condition.
Our proofs in this section make use of his ideas.
\end{rema}

\begin{thebibliography}{CHY21}
\bibitem[Bey75]{beynon} W. M. Beynon, Duality theorems for finitely generated vector lattices, Proc. London Math. Soc. (3) 31 (1975), 114--128.
\bibitem[CHY21]{BP} P. Cubides Kovacsics, M. Hils, and J. Ye, Beautiful pairs, 2021, arXiv:2112.00651.
\bibitem[Ele18]{elef} P. Eleftheriou, Semilinear stars are contractible, Fund. Math. 241 (2018), no. 3, 291--312.
\bibitem[Gla99]{glass} A. M. W. Glass, Partially ordered groups, Series in Algebra, vol. 7, World Scientific Publishing Co., Inc., River Edge, NJ, 1999.
\bibitem[Ovc02]{Ovchinnikov} S. Ovchinnikov, Max-min representation of piecewise linear functions, Beitr. Algebra Geom. 43 (2002), no. 1, 297--302.
\end{thebibliography}
\end{document}
